\begin{textsample}{POS Dim 2 – ba – Score 28.00 – ba/ba\_em\_4.txt}  \label{ex:f2_pos_036}
Tabela 1.
Normativos e \textbf{legislações} das modalidades de ensino | Modalidade | Normativos | Finalidades | |---|---|---| | Educação Especial | Resolução CNE/CEB nº 04, de 2 de outubro de 2009 | \textbf{Institui} as \textbf{Diretrizes} Operacionais para o Atendimento Educacional Especializado ( AEE ) na Educação Básica. | | Educação Especial | Decreto Presidencial nº 6.949, de 25 de agosto de 2009 | Ratifica a Convenção sobre os Direitos das Pessoas com Deficiência/ONU.
Resolução nº 04/2009 CNE/CEB \textbf{institui} as \textbf{Diretrizes} Operacionais para o Atendimento Educacional Especializado ( AEE ) na Educação Básica. | | Educação Especial | Resolução CEE/CEB nº 79, de 15 de setembro de 2009 | Estabelece normas para a Educação Especial na Perspectiva da Educação Inclusiva para todas as etapas e modalidades da Educação Básica no Sistema Estadual de Ensino da Bahia. | | Educação Especial | Nota \textbf{Técnica} SEESP/GAB nº 11, de 7 de maio de 2010 | Dispõe sobre orientações para a institucionalização da \textbf{oferta} do Atendimento Educacional Especializado ( AEE ) em Salas de Recursos Multifuncionais ( SRM ) implantadas nas escolas regulares. | | Educação Especial | Decreto Presidencial nº 7.611, de 17 de novembro de 2011 | Dispõe sobre Educação Especial e Atendimento Educacional Especializado ( AEE ) e dá outras \textbf{providências}. \textbf{Revoga} o Decreto nº 6.571/2008.
Dispõe sobre a classe especial nas escolas regulares e escolas especiais e \textbf{fortalecimento} das instituições especializadas. | | Educação Especial | Lei Brasileira de Inclusão nº 13.146, de 6 de julho de 2015 | \textbf{Institui} a Lei Brasileira de Inclusão da Pessoa com Deficiência ( Estatuto da Pessoa com Deficiência ). | | Educação Especial | Documento \textbf{Orientador}, de 10 de outubro de 2017 | \textbf{Diretrizes} da Educação Inclusiva no Estado da Bahia. | | Educação de Jovens e Adultos | \textbf{Portaria} SEC nº 5.136, de 23 de junho de 2011 | Estabelece normas sobre o procedimento de \textbf{certificação} da \textbf{escolaridade} de jovens e adultos no nível de conclusão do Ensino Fundamental e Médio, por meio dos resultados obtidos no Exame Nacional para \textbf{Certificação} de Competências de Jovens e Adultos ( ENCCEJA ) e Exame Nacional do Ensino Médio ( ENEM ). | | Educação de Jovens e Adultos | Resolução CNE/CEB nº 3, de 15 de junho de 2010 | \textbf{Institui} \textbf{Diretrizes} Operacionais para a Educação de Jovens e Adultos nos aspectos relativos à duração dos \textbf{cursos} e idade mínima para ingresso nos \textbf{cursos} de EJA, idade mínima e \textbf{certificação} nos exames de EJA e Educação de Jovens e Adultos desenvolvida por meio da Educação a Distância. | | Educação de Jovens e Adultos | Resolução CEE-BA nº 239, de 12 de dezembro de 2011 | Dispõe sobre a \textbf{oferta} de Educação de Jovens e Adultos no Estado da Bahia. | | Educação do Campo | Resolução CNE/CEE nº 2, de 28 de abril de 2008 | \textbf{Institui} as \textbf{Diretrizes} Operacionais para a Educação Básica nas Escolas do Campo, um conjunto de princípios e procedimentos para serem observados nos projetos das instituições que integram os diversos sistemas de ensino. | | Educação do Campo | Parecer CNE/CEB nº 1, de 18 de abril de 2006 | Recomenda a adoção da Pedagogia da \textbf{Alternância} em escolas do campo. | | Educação do Campo | Resolução CEE/CEB nº 2, de 28 de abril de 2008 | Estabelece \textbf{diretrizes} complementares, normas e princípios para o desenvolvimento de \textbf{políticas} públicas de atendimento da Educação Básica do Campo. | | Educação do Campo | Decreto Presidencial nº 7.352, de 4 de novembro de 2010 | Dispõe sobre a Política Nacional de Educação do Campo e o Programa Nacional de Educação na \textbf{Reforma} Agrária ( PRONERA ). | | Educação do Campo | Resolução CEE nº 103, de 28 de setembro de 2015 | Dispõe sobre a \textbf{oferta} da Educação do Campo no Sistema Estadual de Ensino da Bahia. | | Educação do Campo | Lei Estadual nº 11.352, de 23 de dezembro de 2008 | \textbf{Institui} o Programa Estadual de Apoio Técnico-Financeiro às Escolas Família Agrícola ( EFAs ) e às Escolas Familiares Rurais ( EFRs ) do Estado da Bahia. | | Educação Escolar Indígena | Resolução CNE/CEB nº 3, de 10 de novembro de 1999 | Fixa Diretrizes Nacionais para o funcionamento das escolas indígenas e dá outras \textbf{providências}. | | Educação Escolar Indígena | Resolução CNE/CEB nº 13, de 10 de maio de 2012 | Define as \textbf{Diretrizes} Curriculares Nacionais para a Educação Escolar Indígena na Educação Básica. | | Educação Escolar Indígena | Lei Nacional nº 11.645, de 10 de março de 2008 | Inclui no \textbf{currículo} oficial da \textbf{rede} de ensino a obrigatoriedade da temática História e Cultura Afro-Brasileira e Indígena. | | Educação Escolar Indígena | \textbf{Portaria} SEC nº 3.918, de 14 de abril de 2012 | Dispõe sobre a reorganização curricular das unidades escolares da educação escolar indígena integrantes da \textbf{rede} pública \textbf{estadual}. | | Educação Escolar Quilombola | Resolução CNE/CEB nº 8, de 20 de novembro de 2012 | Fixa as \textbf{Diretrizes} Curriculares Nacionais para a Educação Escolar Quilombola. | | Educação do Campo, Educação Escolar Indígena e Educação Escolar Quilombola | Lei nº 12.960, de 27 de março de 2014 | Altera a LDBEN para constar a exigência de manifestação de \textbf{órgão} normativo do sistema de ensino ( conselho ) para o fechamento de escolas do campo, indígenas e quilombolas, considerando para tanto a justificativa apresentada pela Secretaria de Educação, a análise do diagnóstico do impacto da ação e a manifestação da comunidade escolar. | | Educação do Campo | Decreto Governamental nº 14.110, de 28 de agosto de 2012 | \textbf{Institui} o Programa Estadual de Apoio Técnico-Financeiro às Escolas Família Agrícola ( EFAs ) e às Escolas Familiares Rurais ( EFRs ) do Estado da Bahia, através de \textbf{entidades} sem fins lucrativos [... ]. | | Educação Escolar Quilombola | Resolução CEE nº 68, de 30 de julho de 2013 | Estabelece normas complementares para \textbf{implantação} e funcionamento das \textbf{Diretrizes} Curriculares Nacionais para a Educação Escolar Quilombola na Educação Básica, no Sistema Estadual de Ensino da Bahia. | Ao longo do documento, outros normativos e \textbf{legislações} serão abordados por terem maior relação com os \textbf{capítulos} tratados no Documento.
Dessa forma, situará melhor o/a leitor/a dentro do contexto dos \textbf{capítulos}.
A seguir, serão reapresentados os normativos e \textbf{legislações} que embasam os temas integradores do DCRB.
A tabela a seguir foi retirada do Volume 1 e será reapresentada mais uma vez, pois deverão ser trabalhados de forma transversal na organização curricular das Unidades Escolares da Bahia, bem como outros temas de relevância para os contextos territoriais, regionais e locais.
Tabela 2.
Normativos e \textbf{legislações} dos temas integradores da Educação Básica | Temas integradores | Normativos | Finalidades | |---|---|---| | Educação para a diversidade ; Educação para as relações de gênero e sexualidade | Lei Federal nº 11.340, de 7 de agosto de 2006 | Cria mecanismos para coibir a violência doméstica e familiar contra a mulher, nos termos do § 8º do art. 226 da Constituição Federal, da Convenção sobre a Eliminação de Todas as Formas de Discriminação contra as Mulheres e da Convenção Interamericana para Prevenir, Punir e Erradicar a Violência contra a Mulher ; dispõe sobre a criação dos Juizados de Violência Doméstica e Familiar contra a Mulher ; altera o Código de Processo Penal, o Código Penal e a Lei de Execução Penal e dá outras \textbf{providências}. | | Educação para as relações de gênero e sexualidade | Lei Federal nº 2.848, § 7º ao art. 121 do Código Penal, de 7 de dezembro de 1940 | Estabelece o aumento da pena do feminicídio. | | Educação para as relações de gênero e sexualidade | Plano Nacional de Políticas para as Mulheres ( 2013–2015 ) | Objetiva o \textbf{fortalecimento} e a institucionalização da Política Nacional para as Mulheres. | | Educação para as relações de gênero e sexualidade | Resolução CEE-BA nº 45, de 24 de agosto de 2020 | Dispõe sobre educação das relações de gêneros e sexualidades no Sistema Estadual de Ensino da Bahia. | | Educação para as relações de gênero e sexualidade | Resolução CEE nº 120, de 5 de novembro de 2013 | Dispõe sobre a inclusão do nome social dos/das estudantes travestis, transexuais e outros no tratamento nos registros escolares e acadêmicos nas Instituições de Ensino que integram o Sistema de Ensino do Estado da Bahia e dá outras \textbf{providências}. | | Educação para as relações de gênero e sexualidade | Lei Federal nº 14.164, de 10 de junho de 2021 | Altera a Lei nº 9.394, de 20 de dezembro de 1996 ( Lei de \textbf{Diretrizes} e Bases da Educação Nacional ), para incluir conteúdo sobre a prevenção da violência contra a mulher nos \textbf{currículos} da educação básica, e \textbf{institui} a Semana Escolar de Combate à Violência contra a Mulher. | | Educação das relações étnico-raciais | Lei nº 10.639, de 9 de janeiro de 2003 | Inclui no \textbf{currículo} oficial da Rede de Ensino a obrigatoriedade da temática História e Cultura Afro-Brasileira e dá outras \textbf{providências}. | | Educação das relações étnico-raciais | Resolução CNE/CEB nº 1, de 17 de junho de 2004 | Estabelece as \textbf{Diretrizes} Curriculares Nacionais para a Educação das Relações Étnico-Raciais e para o Ensino de História e Cultura Afro-Brasileira e Africana. | | Educação das relações étnico-raciais | Lei Federal nº 11.645, de 10 de março de 2008 | Estabelece as \textbf{diretrizes} e bases da educação nacional, para incluir no \textbf{currículo} oficial da \textbf{rede} de ensino a obrigatoriedade da temática História e Cultura Afro-Brasileira e Indígena. | | Educação das relações étnico-raciais | Lei Federal nº 12.288, de 20 de julho de 2010 | \textbf{Institui} o Estatuto da Igualdade Racial ; altera as Leis nº 7.716, de 5 de janeiro de 1989 ; nº 9.029, de 13 de abril de 1995 ; nº 7.347, de 24 de julho de 1985 ; e nº 10.778, de 24 de novembro de 2003. | | Educação das relações étnico-raciais | Lei Federal nº 12.391, de 4 de março de 2011 | Inscreve no Livro dos Heróis da Pátria os nomes dos heróis da “ Revolta dos Búzios ” João de Deus do Nascimento, Lucas Dantas de Amorim Torres, Manuel Faustino Santos Lira e Luís Gonzaga das Virgens e Veiga. | | Educação em Direitos Humanos | Lei Estadual nº 13.182, de 6 de junho de 2014 | \textbf{Institui} o Estatuto da Igualdade Racial e de Combate à Intolerância Religiosa do Estado da Bahia. | | Educação em Direitos Humanos | Lei Estadual nº 13.693, de 11 de janeiro de 2017 | \textbf{Institui} o Dia Estadual de Comemoração ao Dia da África no âmbito do Estado da Bahia. | | Educação em Direitos Humanos | Decreto Estadual nº 16.320, de 21 de setembro de 2015 | \textbf{Institui}, no âmbito do Estado da Bahia, a Década Estadual Afrodescendente, com início em 1º de janeiro de 2015 e final em 31 de dezembro de 2024. | | Educação em Direitos Humanos | Decreto Presidencial nº 7.037, de 21 de dezembro de 2009 | Aprova o Programa Nacional de Direitos Humanos ( PNDH-3 ) e dá outras \textbf{providências}. | | Educação em Direitos Humanos | Plano Nacional de Educação em Direitos Humanos/2007 | Difunde a cultura de Direitos Humanos no país. | | Educação em Direitos Humanos | Decreto Governamental nº 12.019, de 22 de março de 2010 | Aprova o Plano Estadual de Direitos Humanos da Bahia ( PEDH ) e dá outras \textbf{providências}. | | Educação em Direitos Humanos | Plano Estadual de Educação em Direitos Humanos/2009 | Expressa o compromisso do Governo do Estado da Bahia com a promoção da cidadania e dos Direitos Humanos. | | Educação em Direitos Humanos | Parecer CEE/CEB nº 8, de 30 de maio de 2012 | Estabelece \textbf{Diretrizes} Nacionais para a Educação em Direitos Humanos. | | Educação em Direitos Humanos | Lei nº 10.741, de 1º de abril de 2003 | \textbf{Destinado} a regular os direitos assegurados às pessoas com idade igual ou superior a 60 ( \textbf{sessenta} ) anos. ( Estatuto do Idoso ) | | Educação para a diversidade | Parecer CNE/CEB nº 14, de 7 de dezembro de 2011 | \textbf{Diretrizes} para o atendimento de educação escolar de crianças, \textbf{adolescentes} e jovens em situação de itinerância. | | Educação para a diversidade | Resolução CNE/CEB nº 3, de 16 de maio de 2012 | Define \textbf{diretrizes} para o atendimento de educação escolar para populações em situação de itinerância. | | Educação para as relações de gênero e sexualidade | Resolução CEE-BA nº 45, de 24 de agosto de 2020 | Dispõe sobre educação das relações de gêneros e sexualidades no Sistema Estadual de Ensino da Bahia. | | Educação Ambiental | Resolução CNE/CP nº 2, de 15 de junho de 2012 | Estabelece as \textbf{Diretrizes} Curriculares Nacionais para a Educação Ambiental. | | Educação Ambiental | Resolução CEE nº 11, de 17 de janeiro de 2017 | Dispõe sobre a Educação Ambiental no Sistema Estadual de Ensino da Bahia. | | Educação Ambiental | Lei Estadual nº 12.056, de 7 de janeiro de 2011 | \textbf{Institui} a Política de Educação Ambiental do Estado da Bahia e dá outras \textbf{providências}. | | Educação Ambiental | Decreto nº 19.083, de 6 de junho de 2019 | Regulamenta a Lei nº 12.056, de 7 de janeiro de 2011, que \textbf{institui} a Política de Educação Ambiental do Estado da Bahia e dá outras \textbf{providências}. | | Educação Ambiental | Lei Estadual nº 13.572, de 30 de agosto de 2016 | \textbf{Institui} a Política Estadual de Convivência com o Semiárido e o Sistema Estadual de Convivência com o Semiárido e dá outras \textbf{providências}. | | Saúde na Escola | Decreto Presidencial nº 6.286, de 5 de dezembro de 2007 | \textbf{Institui} o Programa Saúde na Escola. | | Saúde na Escola | \textbf{Portaria} nº 2.728, de 7 de abril de 2016 | \textbf{Institui} a Promoção da Saúde e Prevenção de Doenças e Agravos no contexto escolar, com ênfase no combate ao mosquito Aedes aegypti. | | Saúde na Escola | \textbf{Portaria} Conjunta SEPLAN/SESAB/SEC nº 001, de 17 de maio de 2014 | Institucionaliza as ações transversais e esforços \textbf{intersetoriais} para \textbf{implantação} do Programa de Ação Estadual de Prevenção da gravidez e assistência ao parto na adolescência. | | Saúde na Escola | Lei Federal nº 11.947, de 26 de junho de 2009 | Dispõe sobre o atendimento da alimentação escolar e do Programa Dinheiro Direto na Escola aos/às estudantes da Educação Básica ; altera as Leis nº 10.880, de 9 de junho de 2004 ; nº 11.273, de 6 de fevereiro de 2006 ; nº 11.507, de 20 de julho de 2007 ; e \textbf{revoga} dispositivos da Medida Provisória nº 2.178-36, de 24 de agosto de 2001, e a Lei nº 8.913, de 12 de julho de 1994, e dá outras \textbf{providências}. | | Saúde na Escola | \textbf{Portaria} Conjunta SESAB/SEC nº 01, de 29 de agosto de 2018 | Dispõe sobre a obrigatoriedade da apresentação da carteira/cartão de vacinação em creches e escolas, em todo o território do Estado da Bahia. | | Saúde na Escola | Lei Federal nº 13.666, de 16 de maio de 2018 | Altera a Lei nº 9.394, de 20 de dezembro de 1996 ( Lei de \textbf{Diretrizes} e Bases da Educação Nacional ), para incluir o tema transversal da educação alimentar e nutricional no \textbf{currículo} escolar. | | Educação Fiscal | \textbf{Portaria} Interministerial MF/MEC nº 413, de 31 de dezembro de 2002 | Implementa o Programa Nacional de Educação Fiscal. | | Educação Fiscal | Decreto Estadual nº 15.737, de 10 de dezembro de 2014 | \textbf{Institui} a Educação Fiscal na Bahia. | Fonte : Elaboração Equipe SEC/SUDEP ( 2018 ).

% matched lemmas: adolescente, alternância, capítulo, certificação, currículo, curso, destinar, diretriz, entidade, escolaridade, estadual, fortalecimento, implantação, instituir, intersetorial, legislação, oferta, orientador, política, portaria, providência, rede, reforma, revogar, sessenta, técnico, órgão
\end{textsample}
